% ═══════════════════════════════════════════════════════════════
% ARBEITSBLATT: Dichte (Klasse 7) - KOMPAKT (2 Seiten, 2-spaltig)
% Begleitend zum digitalen Lernpfad
% ═══════════════════════════════════════════════════════════════

\documentclass[10pt, a4paper]{article}

% ═══════════════════════════════════════════════════════════════
% PAKETE
% ═══════════════════════════════════════════════════════════════

\usepackage[utf8]{inputenc}
\usepackage[T1]{fontenc}
\usepackage[ngerman]{babel}
\usepackage{helvet}
\renewcommand{\familydefault}{\sfdefault}

\usepackage[margin=1.5cm, top=1.8cm, bottom=1.5cm]{geometry}
\usepackage{amsmath, amssymb}
\usepackage{siunitx}
\sisetup{locale = DE, per-mode = symbol}

\usepackage{graphicx}
\usepackage{tikz}
\usepackage{enumitem}
\usepackage{tabularx}
\usepackage{booktabs}
\usepackage{array}
\usepackage{multicol}
\usepackage{setspace}

\usepackage{xcolor}
\usepackage{tcolorbox}
\usepackage{fancyhdr}
\usepackage{lastpage}
\usepackage{qrcode}

% Kompaktere Listen
\setlist{nosep, leftmargin=1.5em}

% ═══════════════════════════════════════════════════════════════
% FARBEN
% ═══════════════════════════════════════════════════════════════

\definecolor{physik}{HTML}{2c5aa0}
\definecolor{grau}{HTML}{888888}
\definecolor{hellgrau}{HTML}{f0f0f0}
\definecolor{success}{HTML}{2a7a4b}

\newcommand{\fachfarbe}{physik}

% ═══════════════════════════════════════════════════════════════
% VARIABLEN
% ═══════════════════════════════════════════════════════════════

\newcommand{\thema}{Dichte}
\newcommand{\fach}{Physik}
\newcommand{\klasse}{7}

% ═══════════════════════════════════════════════════════════════
% KOPF- UND FUSSZEILE
% ═══════════════════════════════════════════════════════════════

\pagestyle{fancy}
\fancyhf{}
\renewcommand{\headrulewidth}{0.4pt}
\renewcommand{\footrulewidth}{0pt}

\lhead{\small\fach{} Kl.\,\klasse{} -- \thema}
\rhead{\small Name: \rule{4cm}{0.4pt}}

\cfoot{\small\thepage/\pageref{LastPage}}

% ═══════════════════════════════════════════════════════════════
% BEFEHLE
% ═══════════════════════════════════════════════════════════════

% Lücke
\newcommand{\luecke}[1]{\rule{#1}{0.4pt}}

% Kompakte Formelbox
\newtcolorbox{formelbox}{
    colback=hellgrau,
    colframe=\fachfarbe,
    boxrule=0.8pt,
    arc=0pt,
    left=6pt, right=6pt, top=6pt, bottom=6pt
}

% Kompakter Merksatz
\newtcolorbox{merksatz}{
    colback=hellgrau,
    colframe=success,
    boxrule=0.8pt,
    arc=0pt,
    left=6pt, right=6pt, top=4pt, bottom=4pt,
    fonttitle=\bfseries\small,
    title=Merksatz
}

% Abschnittstitel
\newcommand{\abschnitt}[1]{%
    \vspace{6pt}
    \noindent\textbf{\textcolor{\fachfarbe}{#1}}
    \vspace{3pt}
}

% ═══════════════════════════════════════════════════════════════
% DOKUMENT
% ═══════════════════════════════════════════════════════════════

\begin{document}
\setstretch{1.1}

% ═══════════════════════════════════════════════════════════════
% TITEL
% ═══════════════════════════════════════════════════════════════

\begin{center}
{\large\bfseries Arbeitsblatt: \thema}\\[2pt]
{\small\textcolor{grau}{Begleitend zum Lernpfad -- in den Hefter übernehmen}}
\end{center}

\vspace{-2pt}
\hrule
\vspace{8pt}

% ═══════════════════════════════════════════════════════════════
% ZWEISPALTIG
% ═══════════════════════════════════════════════════════════════

\begin{multicols}{2}
\small

% ─────────────────────────────────────────────────────────────
% 1. DEFINITION
% ─────────────────────────────────────────────────────────────

\abschnitt{1. Was ist Dichte?}

\textbf{Definition:} Die Dichte gibt an, wie viel

\vspace{3pt}
\luecke{2.2cm} in einem bestimmten \luecke{2.2cm}

\vspace{3pt}
\noindent eines Stoffes enthalten ist.

\vspace{5pt}
\noindent Die Dichte ist eine \luecke{3cm}.

\vspace{2pt}
{\scriptsize\textcolor{grau}{(Masse, Volumen, Stoffeigenschaft)}}

\vspace{8pt}

% ─────────────────────────────────────────────────────────────
% 2. FORMEL
% ─────────────────────────────────────────────────────────────

\abschnitt{2. Die Formel}

\begin{formelbox}
\centering
{\large $\rho = \dfrac{m}{V}$}

\vspace{6pt}

\begin{tabular}{r@{\,=\,}l@{\quad}l}
$\rho$ & \luecke{1.8cm} & $\si{g/cm^3}$ \\[4pt]
$m$ & \luecke{1.8cm} & $\si{g}$ \\[4pt]
$V$ & \luecke{1.8cm} & $\si{cm^3}$
\end{tabular}
\end{formelbox}

\vspace{4pt}

\noindent\textbf{Umstellung:}

\vspace{2pt}

\begin{tabular}{|l|l|}
\hline
Gesucht & Formel \\
\hline
$\rho$ & $\rho = m / V$ \\
\hline
$m$ & $m = $ \luecke{1.5cm} \\
\hline
$V$ & $V = $ \luecke{1.5cm} \\
\hline
\end{tabular}

\vspace{8pt}

% ─────────────────────────────────────────────────────────────
% 3. DICHTEWERTE
% ─────────────────────────────────────────────────────────────

\abschnitt{3. Typische Dichtewerte}

\begin{tabular}{|l|r|c|}
\hline
\textbf{Stoff} & $\rho$ \scriptsize{(g/cm³)} & \textbf{in Wasser} \\
\hline
Kork & 0,2 & \luecke{1cm} \\
\hline
Holz & 0,5 & \luecke{1cm} \\
\hline
Öl & 0,9 & \luecke{1cm} \\
\hline
Eis & 0,92 & \luecke{1cm} \\
\hline
\textbf{Wasser} & \textbf{1,0} & --- \\
\hline
Plastik & 1,05 & \luecke{1cm} \\
\hline
Alu & 2,7 & \luecke{1cm} \\
\hline
Eisen & 7,9 & \luecke{1cm} \\
\hline
Blei & 11,3 & \luecke{1cm} \\
\hline
\end{tabular}

{\scriptsize\textcolor{grau}{(schwimmt / sinkt)}}

\vspace{8pt}

% ─────────────────────────────────────────────────────────────
% 4. SCHWIMMEN/SINKEN
% ─────────────────────────────────────────────────────────────

\abschnitt{4. Schwimmen, Schweben, Sinken}

\begin{merksatz}
\begin{tabular}{|l|l|}
\hline
\textbf{Bedingung} & \textbf{Verhalten} \\
\hline
$\rho_\text{K}$ \luecke{0.6cm} $\rho_\text{Fl}$ & Schwimmen \\
\hline
$\rho_\text{K}$ \luecke{0.6cm} $\rho_\text{Fl}$ & Schweben \\
\hline
$\rho_\text{K}$ \luecke{0.6cm} $\rho_\text{Fl}$ & Sinken \\
\hline
\end{tabular}

\vspace{2pt}
{\scriptsize (Ergänze: $<$ , $=$ , $>$)}
\end{merksatz}

% ═══════════════════════════════════════════════════════════
% SPALTE 2
% ═══════════════════════════════════════════════════════════

\columnbreak

% ─────────────────────────────────────────────────────────────
% 5. BEISPIELRECHNUNG
% ─────────────────────────────────────────────────────────────

\abschnitt{5. Beispielrechnung}

\textbf{Aufgabe:} Ein Block hat $m = \SI{216}{g}$ und $V = \SI{80}{cm^3}$. Berechne $\rho$.

\vspace{6pt}

\begin{tabular}{@{}l@{\,}l}
\textbf{Geg.:} & $m = $ \luecke{2cm} \\[4pt]
& $V = $ \luecke{2cm} \\[6pt]
\textbf{Ges.:} & \luecke{2cm} \\[6pt]
\textbf{Formel:} & $\rho = $ \luecke{2cm} \\[6pt]
\textbf{Eins.:} & $\rho = \dfrac{216\,\text{g}}{80\,\text{cm}^3} = $ \luecke{1.8cm} \\[8pt]
\textbf{Antw.:} & Der Block hat die Dichte \\[2pt]
& \luecke{2.5cm}. \\[2pt]
& Das ist \luecke{2cm}.
\end{tabular}

\vspace{10pt}

% ─────────────────────────────────────────────────────────────
% 6. SIMULATION
% ─────────────────────────────────────────────────────────────

\abschnitt{6. Beobachtungen (Simulation)}

\textbf{a) In Wasser} ($\rho = \SI{1,0}{g/cm^3}$):

\vspace{3pt}

\begin{tabular}{|l|c|c|c|}
\hline
Material & schw. & schwe. & sinkt \\
\hline
Kork (0,2) & $\square$ & $\square$ & $\square$ \\
\hline
Holz (0,5) & $\square$ & $\square$ & $\square$ \\
\hline
Öl (0,9) & $\square$ & $\square$ & $\square$ \\
\hline
Eis (0,92) & $\square$ & $\square$ & $\square$ \\
\hline
Fisch (1,0) & $\square$ & $\square$ & $\square$ \\
\hline
Eisen (7,9) & $\square$ & $\square$ & $\square$ \\
\hline
Eigener: \luecke{0.5cm} & $\square$ & $\square$ & $\square$ \\
\hline
\end{tabular}

\vspace{6pt}

\textbf{b) Vergleich: Salzwasser} ($\rho = \SI{1,03}{g/cm^3}$)

\vspace{3pt}

Was passiert mit Plastik ($\rho = 1,05$)?

\vspace{3pt}

\begin{tabular}{@{}l@{ }l}
Wasser: & \luecke{3.5cm} \\[4pt]
Salzw.: & \luecke{3.5cm}
\end{tabular}

\vspace{4pt}

\textbf{Erklärung:} \luecke{4cm}

\vspace{3pt}
\luecke{5cm}

\vspace{10pt}

% ─────────────────────────────────────────────────────────────
% 7. ZUSAMMENFASSUNG
% ─────────────────────────────────────────────────────────────

\abschnitt{7. Zusammenfassung}

\begin{merksatz}
Die \textbf{Dichte} $\rho$ gibt an, wie viel Masse in einem Volumen steckt.

\vspace{3pt}

$\rho = \dfrac{m}{V}$ \quad mit \quad $[\rho] = \SI{}{g/cm^3}$

\vspace{3pt}

$\rho_\text{K} < \rho_\text{Fl}$ $\rightarrow$ \textbf{schwimmt}

$\rho_\text{K} = \rho_\text{Fl}$ $\rightarrow$ \textbf{schwebt}

$\rho_\text{K} > \rho_\text{Fl}$ $\rightarrow$ \textbf{sinkt}
\end{merksatz}

\vspace{8pt}

% ─────────────────────────────────────────────────────────────
% QR-CODE
% ─────────────────────────────────────────────────────────────

\begin{center}
\fbox{%
\begin{minipage}{0.85\linewidth}
\centering
\vspace{4pt}
\textbf{\small Digitaler Lernpfad}

\vspace{4pt}

% URL vor Nutzung anpassen:
\qrcode[height=1.5cm]{https://example.com/lernpfad/kl7-dichte}

\vspace{2pt}
{\tiny\textcolor{grau}{QR-Code scannen}}
\vspace{2pt}
\end{minipage}%
}
\end{center}

\end{multicols}

% ═══════════════════════════════════════════════════════════════
% SEITE 2: ÜBUNGSAUFGABEN
% ═══════════════════════════════════════════════════════════════

\newpage

\begin{center}
{\large\bfseries Übungsaufgaben: \thema}
\end{center}

\vspace{-2pt}
\hrule
\vspace{8pt}

\begin{multicols}{2}
\small

% ─────────────────────────────────────────────────────────────
% NIVEAU 1
% ─────────────────────────────────────────────────────────────

\abschnitt{Niveau (*) -- Basis}

\textbf{A1.} Berechne die Dichte.

\vspace{3pt}

\begin{tabular}{|c|c|c|}
\hline
$m$ & $V$ & $\rho$ \\
\hline
\SI{200}{g} & \SI{100}{cm^3} & \luecke{1.5cm} \\
\hline
\SI{500}{g} & \SI{250}{cm^3} & \luecke{1.5cm} \\
\hline
\SI{810}{g} & \SI{300}{cm^3} & \luecke{1.5cm} \\
\hline
\end{tabular}

\vspace{8pt}

\textbf{A2.} Schwimmt oder sinkt der Körper in Wasser?

\vspace{3pt}

\begin{tabular}{|l|c|}
\hline
Stoff mit $\rho$ & Verhalten \\
\hline
Kork ($\rho = 0,2$) & \luecke{1.8cm} \\
\hline
Stein ($\rho = 2,5$) & \luecke{1.8cm} \\
\hline
Wachs ($\rho = 0,9$) & \luecke{1.8cm} \\
\hline
Glas ($\rho = 2,5$) & \luecke{1.8cm} \\
\hline
\end{tabular}

\vspace{8pt}

\textbf{A3.} Ergänze die Lücken:

\vspace{3pt}

a) Ein Körper schwimmt, wenn seine Dichte

\luecke{1.5cm} als die der Flüssigkeit ist.

\vspace{4pt}

b) Ein Körper sinkt, wenn seine Dichte

\luecke{1.5cm} als die der Flüssigkeit ist.

\vspace{4pt}

c) Ein Körper schwebt, wenn beide Dichten

\luecke{1.5cm} sind.

\vspace{10pt}

% ─────────────────────────────────────────────────────────────
% NIVEAU 2
% ─────────────────────────────────────────────────────────────

\abschnitt{Niveau (**) -- Erweiterung}

\textbf{A4.} Berechne die fehlende Größe.

\vspace{3pt}

a) $\rho = \SI{8}{g/cm^3}$, $V = \SI{25}{cm^3}$

\vspace{2pt}
$m = $ \luecke{3cm}

\vspace{5pt}

b) $\rho = \SI{2,7}{g/cm^3}$, $m = \SI{540}{g}$

\vspace{2pt}
$V = $ \luecke{3cm}

\vspace{8pt}

\textbf{A5.} Ein Würfel hat die Kantenlänge $a = \SI{4}{cm}$ und die Masse $m = \SI{128}{g}$.

\vspace{3pt}

a) Berechne das Volumen: $V = a^3 = $ \luecke{2cm}

\vspace{4pt}

b) Berechne die Dichte: $\rho = $ \luecke{2cm}

\vspace{4pt}

c) Schwimmt der Würfel? \luecke{2.5cm}

% ═══════════════════════════════════════════════════════════
% SPALTE 2 SEITE 2
% ═══════════════════════════════════════════════════════════

\columnbreak

% ─────────────────────────────────────────────────────────────
% NIVEAU 3
% ─────────────────────────────────────────────────────────────

\abschnitt{Niveau (***) -- Gymnasium}

\textbf{A6.} Ein Metallstück verdrängt in Wasser \SI{50}{mL} und wiegt \SI{395}{g}.

\vspace{3pt}

a) Volumen: $V = $ \luecke{2cm}

\vspace{4pt}

b) Dichte: $\rho = $ \luecke{2cm}

\vspace{4pt}

c) Um welches Metall handelt es sich?

\vspace{2pt}
\luecke{4cm}

\vspace{8pt}

\textbf{A7.} Du mischst \SI{100}{mL} Wasser ($\rho = 1,0$) mit \SI{50}{mL} Öl ($\rho = 0,8$).

\vspace{3pt}

a) $m_\text{Wasser} = $ \luecke{2cm}

\vspace{4pt}

b) $m_\text{Öl} = $ \luecke{2cm}

\vspace{4pt}

c) $m_\text{gesamt} = $ \luecke{2cm}

\vspace{4pt}

d) Welche Schicht ist oben? \luecke{2cm}

\vspace{10pt}

\textbf{A8.} Ein U-Boot kann abtauchen und auftauchen.

\vspace{3pt}

a) Was macht das U-Boot mit seiner Dichte, um abzutauchen?

\vspace{2pt}
\luecke{5cm}

\vspace{5pt}

b) Wie erreicht es das technisch?

\vspace{2pt}
\luecke{5cm}

\vspace{2pt}
\luecke{5cm}

\vspace{12pt}

% ─────────────────────────────────────────────────────────────
% SELBSTKONTROLLE
% ─────────────────────────────────────────────────────────────

\abschnitt{Selbstkontrolle}

\begin{tabular}{|p{4.5cm}|c|}
\hline
\textbf{Ich kann...} & \\
\hline
...die Definition der Dichte nennen & $\square$ \\
\hline
...die Formel $\rho = m/V$ anwenden & $\square$ \\
\hline
...die Formel umstellen & $\square$ \\
\hline
...vorhersagen, ob ein Körper schwimmt & $\square$ \\
\hline
...erklären, warum Holz schwimmt & $\square$ \\
\hline
\end{tabular}

\vspace{10pt}

% ─────────────────────────────────────────────────────────────
% LÖSUNGEN (klein)
% ─────────────────────────────────────────────────────────────

{\scriptsize\textcolor{grau}{
\textbf{Lösungen:}
A1: 2; 2; 2,7 |
A4a: 200g | A4b: 200cm³ |
A5: 64cm³; 2g/cm³; sinkt |
A6: 50cm³; 7,9g/cm³; Eisen |
A7: 100g; 40g; 140g; Öl
}}

\end{multicols}

\end{document}
